\cchapter{ВВЕДЕНИЕ}                         % Заголовок
\addcontentsline{toc}{chapter}{ВВЕДЕНИЕ}    % Добавляем его в оглавление


Автомобиль сегодня является основным средством передвижения для большой части населения, а ДТП и последующий ремонт - нередкое явление. При повреждении автомобиля, перед тем, как отдать его в ремонт, владельцы хотят иметь представление, сколько потребуется платить за услуги и ожидать их. Однако получить такую информацию быстро и не приезжая в автосервис выходит далеко не всегда.

Сотрудники автосервисов ежедневно вынуждены выделять время на обработку такого рода обращений по оценке повреждений от потенциальных клиентов. Если клиент отправляет фотографию мастеру, он вынужден ждать, пока сотрудник выйдет на связь и даст ответ. Если же требуется приехать в автосервис лично, клиент вынужден тратить дополнительное время на дорогу, оставаясь без понимания стоимости и длительности ремонта. Это потенциально может вызывать дискомфорт и снижать вероятность обращения, если, например, в другом автосервисе процессы устроены более прозрачно. В обоих случаях персоналу нужно отвлекаться от работы, чтобы оценивать машины новых клиентов, а клиенту нужно делать много шагов по взаимодействию с персоналом автосервиса, чтобы получить обратную связь по своему автомобилю.

Ранее этот процесс как-то улучшить не представлялось возможным, так как точная оценка повреждений требовала личного взгляда специалиста. Но сейчас, с развитием методов машинного обучения и компьютерного зрения, это стало возможным - теперь можно провести анализ изображения и параметров автомобиля с помощью искусственного интеллекта с достаточной точностью. Это позволяет определить ответы на самые частые вопросы клиентов - стоимость и длительность работ, и перенести часть работы на автоматизированный сервис.

Таким образом, целью работы является повышение эффективности процесса оценки стоимости и длительности ремонта поврежденных автомобилей. 

Сейчас, при высокой конкуренции между автосервисами скорость обработки обращений является одним из важнейших факторов, влияющих на удержание и привлечение новых клиентов. Автосервис, имеющий интеллектуальную систему, которая способна практически моментально дать ответ клиенту по новому обращению, имеет преимущество перед остальными. Клиенту комфортнее взаимодействовать, ведь он получает обратную связь по главным интересующим вопросам почти сразу, а доверие к сервису возрастает. К тому же, на рынке отсутствуют подобного рода решения, выполняющие автоматическую оценку повреждений по фотографии автомобиля, что делает разработку инновационной и востребованной.
